\begin{proposition}[absolute cotangent direct images]
\label{prop:absolute-cotangent-direct-images}
\Verified\ Put $\mathcal K=\ker\kappa_1$.  Then
\[
 f_*\Omega_X^1\simeq\mathcal O(-2),\qquad
 R^1f_*\Omega_X^1\simeq\mathcal K,\qquad
 R^2f_*\Omega_X^1\simeq\mathcal O\oplus\mathcal O(-1).
\tag{11.13}
\]
None of these sheaves has torsion.
\end{proposition}

\begin{proof}
The three normalized components have
\[
 R^0f_*\nu_*\mathcal O_{\widetilde D}
  =\mathbb C_{p_0}\oplus\mathbb C_{p_1}\oplus\mathbb C_{p_2},
 \quad
 R^1f_*\nu_*\mathcal O_{\widetilde D}
  =\mathbb C_{p_1}\oplus\mathbb C_{p_2},
 \quad R^2=0.
\tag{11.14}
\]
Indeed $\mathrm{dP}_6$ is rational and $S_1,S_2$ are bielliptic, with
$h^1(\mathcal O_{S_j})=1$ and $h^2(\mathcal O_{S_j})=0$.

The long direct-image sequence of (11.2), together with
\cref{prop:log-ks-divisor}, gives
\[
 f_*\mathcal A=\mathcal O(1),\qquad
 0\longrightarrow T\longrightarrow R^1f_*\mathcal A
 \longrightarrow\mathcal K\longrightarrow0,
 \qquad R^2f_*\mathcal A=\mathcal O\oplus\mathcal O(-1),
\tag{11.15}
\]
where $T=\mathbb C_{p_1}\oplus\mathbb C_{p_2}$.
The residue of the base logarithmic form is weighted evaluation:
\[
 \mathcal O(1)\longrightarrow
 \mathbb C_{p_0}\oplus\mathbb C_{p_1}\oplus\mathbb C_{p_2},
 \qquad q\longmapsto(q(p_0),\,3q(p_1),\,4q(p_2)).
\tag{11.16}
\]
It is surjective as a sheaf and has kernel
$\mathcal O(1)(-p_0-p_1-p_2)=\mathcal O(-2)$.

It remains to identify the torsion in (11.15).  At $p_j$, the residue map on
$G\otimes\mathbb C_{p_j}$ sends
\[
 x b_0+y s^ab_1\longmapsto m_jx[b_0]
 \in H^1(S_j,\mathcal O_{S_j}).
\tag{11.17}
\]
More precisely, after trivializing $\mathcal L$ by
$\ell=dt/t$, the logarithmic conormal and residue maps form the exact
commutative diagram
\[
\begin{array}{ccccccccc}
0&\longrightarrow&f^*\mathcal L&\longrightarrow&\mathcal A
 &\longrightarrow&\mathcal E&\longrightarrow&0\\
&&\big\downarrow\,{\scriptstyle\operatorname{res}_B}
&&\big\downarrow\,{\scriptstyle\operatorname{res}_{S_j}}
&&\big\downarrow\,{\scriptstyle 0}&&\\
0&\longrightarrow&i_{j*}\mathcal O_{S_j}&\xrightarrow{\ =\ }
 i_{j*}\mathcal O_{S_j}&\longrightarrow&0&&
\end{array}
\tag{11.16a}
\]
with
\[
 \operatorname{res}_B(q\ell)=m_jq|_{S_j}.
\tag{11.16b}
\]
Applying $Rf_*$ to (11.16a) gives the typed composite
\[
 f_*\mathcal E|_{p_j}\xrightarrow{\ \kappa_0\ }
 (\mathcal L\otimes G)|_{p_j}
 \xrightarrow{\ \operatorname{res}_{1,j}\ }
 H^1(S_j,\mathcal O_{S_j}),
 \qquad \operatorname{res}_{1,j}\kappa_0=0,
\tag{11.16c}
\]
and (11.16b) gives exactly
\[
 \operatorname{res}_{1,j}(xb_0+ys^ab_1)=m_jx[b_0].
\]
Thus its kernel is the one-dimensional line
$\mathbb C\,s^ab_1$.  The simple zero of $\det\kappa_0$ proved in
\cref{prop:log-ks-divisor} implies that $\kappa_0\bmod t$ has rank one;
the zero composite in (11.16c) therefore forces
\[
 \operatorname{im}(\kappa_0\bmod t)
 =\ker\operatorname{res}_{1,j}
 =\mathbb C\,s^ab_1.
\tag{$11.17a_{\mathrm{fin}}$}
\]
Consequently the induced map
$\operatorname{coker}\kappa_0\to H^1(S_j,\mathcal O_{S_j})$ is an
isomorphism at each finite fibre.  Thus each summand
$T_{p_j}\to H^1(S_j,\mathcal O_{S_j})$ is an isomorphism.  In the long
sequence of (11.12), (11.16) first removes the entire degree-zero residue,
and (11.17) then cancels $T$.  What remains is precisely (11.13).
\end{proof}
